\documentclass{article}

\usepackage[preprint]{neurips_2024}

\usepackage[utf8]{inputenc}
\usepackage[T1]{fontenc}
\usepackage{hyperref}
\usepackage{url}
\usepackage{booktabs}
\usepackage{graphicx}
\usepackage{array}
\usepackage{amsfonts}
\usepackage{nicefrac}
\usepackage{microtype}
\usepackage{xcolor}
\makeatletter
\renewcommand{\@notice}{}
\makeatother

\title{Empirical Evaluation of Update Stability in Continuous-Control Policy Optimization}

\author{%
  Hardik Shrestha \\
  Student ID: 251163014 \\
  Course Final Project \\
}

\begin{document}

\maketitle

\begin{abstract}
Policy-gradient methods are attractive for continuous control, but they are also sensitive to update size. This project studies that sensitivity through a controlled benchmark rather than through a new optimizer. Four closely related on-policy methods were compared under a locked implementation and evaluation protocol: A2C, PPO-Clip, PPO-KL, and TRPO. The benchmark used Pendulum-v1, Hopper-v4, Walker2d-v4, and HalfCheetah-v4, with matched architecture, preprocessing, rollout size, evaluation cadence, and seed budget. On the main locomotion tasks, the step-control methods outperformed A2C under the default configuration. TRPO achieved the highest mean final return on all four environments and zero collapses across 20 runs, although the Hopper-v4 variance was large enough that fine-grained ranking there should be interpreted cautiously. PPO-Clip was the strongest first-order baseline. PPO-KL was functionally distinct because target-KL early stopping was active, but it was less reliable on Walker2d-v4. All quantitative claims in the report were recomputed from per-run \texttt{metrics.csv}, \texttt{updates.csv}, and \texttt{run\_status.json} files rather than from known-bad aggregate exports.
\end{abstract}

\section{Introduction}

Policy-gradient reinforcement learning is appealing because it optimizes a stochastic policy directly and handles continuous action spaces naturally. That makes it a natural fit for robotics-style control tasks such as Hopper, Walker2d, and HalfCheetah, where the action is a vector of torques rather than a small discrete decision. Its practical weakness is that performance is tightly coupled to update size. If policy steps are too conservative, learning is slow; if they are too aggressive, training can degrade abruptly or collapse. That tension appears in foundational policy-gradient work and remains central in modern actor-critic systems [1, 2].

The problem studied in this project is therefore not simply ``which algorithm gets the highest return.'' The more useful problem is whether explicit control of policy movement makes on-policy learning more reliable under a matched continuous-control benchmark. Existing techniques do not fully solve this problem in a plug-and-play way. A2C is simple and efficient, but it relies heavily on ordinary learning-rate tuning. TRPO has a strong theoretical motivation, but it is computationally heavier and can be sensitive to implementation details. PPO was designed to be simpler than TRPO, but the label ``PPO'' covers multiple control mechanisms, including clipping, KL penalties, and target-KL early stopping. Comparing these methods without locking the implementation and evaluation protocol can make the result reflect incidental engineering choices rather than the actual update-control idea.

TRPO and PPO are best understood as responses to that step-size problem. TRPO enforces an explicit KL trust region [4]. PPO replaces second-order trust-region machinery with first-order approximations, most commonly ratio clipping or KL-based control [6]. At the same time, the deep RL evaluation literature warns that seed choice, implementation details, and protocol drift can dominate named algorithm differences [7, 8]. For a course project, that means a fair comparison is itself a technical contribution.

The technique developed in this project is a matched update-stability evaluation framework rather than a new optimizer. The intuition is that policy-update methods should be judged not only by final return, but also by whether they prevent unusable training runs under the same architecture, rollout budget, normalization, evaluation cadence, and seed budget. The framework locks those conditions, defines run-level collapse and update-level instability diagnostics, and then compares A2C, PPO-Clip, PPO-KL, and TRPO. The central question is narrow: under shared architecture, preprocessing, rollout budget, and evaluation cadence, does explicit policy-step control improve default-setting performance and robustness relative to a standard actor-critic baseline in continuous control?

\section{Techniques to Tackle the Problem}

\subsection{Related work and framing}

Prior work gives four relevant ideas. First, REINFORCE shows that policy gradients can be estimated directly from sampled returns, but also exposes the variance problem that makes naive policy-gradient learning hard to tune [1]. Actor-critic methods reduce that variance by learning a value baseline, and asynchronous actor-critic work showed that these ideas can scale to neural-network policies [2]. Second, approximately optimal approximate reinforcement learning and generalized advantage estimation both motivate the idea that policy improvement depends on controlling approximation error and using lower-variance advantage estimates [3, 5]. Third, TRPO and PPO are explicit attempts to limit destructive policy movement while retaining useful policy improvement [4, 6]. TRPO does this with a KL-constrained trust-region update; PPO does it with simpler first-order mechanisms such as clipping or KL control. Fourth, evaluation-focused studies show that implementation parity and reporting discipline are necessary if algorithm comparisons are to mean much at all [7, 8]. The benchmark in this project follows that fourth lesson directly.

This related work also explains why the project compares existing methods rather than claiming a brand-new reinforcement-learning algorithm. The research gap is practical and empirical: students and practitioners often know that TRPO and PPO are supposed to be more stable than a plain actor-critic baseline, but the strength of that claim depends on how carefully the comparison is controlled. A sloppy comparison could give one method a better architecture, different preprocessing, different rollout lengths, or a more forgiving failure rule. The contribution here is to remove those avoidable differences and then ask whether the update-control mechanism still matters.

\subsection{Benchmark design}

The project freezes a single comparison protocol across all four methods. Every algorithm uses separate actor and critic MLPs with hidden sizes \texttt{[64, 64]}, \texttt{tanh} activations, a tanh-squashed diagonal-Gaussian policy, orthogonal initialization, actor hidden gain \texttt{sqrt(2)}, actor output gain \texttt{0.01}, critic output gain \texttt{1.0}, and initial log standard deviation \texttt{-0.5}. Observation normalization is enabled for all methods, reward normalization is disabled, advantage normalization is enabled, and GAE uses \texttt{gamma = 0.99} and \texttt{lambda = 0.95}. Value targets bootstrap on time-limit truncation but not on true terminal transitions. All methods use the same rollout batch of 2048 environment steps per update, including TRPO. Evaluation is run at step 0 and then every 10{,}000 environment steps for 10 episodes with \texttt{deterministic=True}, which for this Gaussian policy means using the mean action under the tanh-squashed policy. The same tanh-squash correction is used across methods when logging policy statistics.

The benchmark logs both return and stability outcomes. Run-level robustness is measured by collapse incidence under the pre-registered relative-performance rule. Update-level behavior is measured with an unstable-update diagnostic based on a KL condition plus an immediate post-update performance drop. That diagnostic is asymmetric by design: TRPO and PPO-KL are flagged when \texttt{mean\_kl\_old\_new > 2 x 0.02 = 0.04}, whereas A2C and PPO-Clip use a fixed threshold of \texttt{0.05}. That asymmetry must be disclosed when interpreting the metric.

\begin{table}[t]
  \caption{Locked benchmark protocol used for the comparison.}
  \label{tab:protocol}
  \centering
  \small
  \begin{tabular}{>{\raggedright\arraybackslash}p{0.31\linewidth}>{\raggedright\arraybackslash}p{0.61\linewidth}}
    \toprule
    Design choice & Value used in the final benchmark \\
    \midrule
    Algorithms & A2C, PPO-Clip, PPO-KL, TRPO \\
    Environments & Pendulum-v1, Hopper-v4, Walker2d-v4, HalfCheetah-v4 \\
    Policy/value network & Separate actor and critic MLPs, hidden sizes \texttt{[64,64]}, \texttt{tanh} activations \\
    Rollout and advantage estimation & 2048 environment steps per update, GAE with \texttt{gamma=0.99} and \texttt{lambda=0.95} \\
    Normalization & Observation normalization enabled, reward normalization disabled, advantage normalization enabled \\
    Evaluation & Step 0 and every 10,000 environment steps, 10 deterministic episodes per checkpoint \\
    Robustness reporting & Collapse incidence as the primary run-level robustness metric, with collapsed runs retained in averages \\
    \bottomrule
  \end{tabular}
\end{table}

\subsection{Compared methods}

A2C is the baseline without explicit trust-region or proximal step control beyond ordinary learning-rate choice. PPO-Clip uses the clipped probability-ratio surrogate, limiting the incentive to move the new policy too far from the old one on sampled actions without solving an exact constrained optimization problem. PPO-KL uses the same rollout, minibatching, actor-critic architecture, and multi-epoch scaffold as PPO-Clip, but its policy loss is the unclipped importance-ratio surrogate \texttt{-(A * ratio)}, and policy epochs stop early once the mean KL exceeds the target. Because SB3 requires a \texttt{clip\_range} argument, the implementation passes a non-operative bound (\texttt{1e9}) while relying on target-KL early stopping as the actual control mechanism. TRPO uses an explicit KL constraint with \texttt{max\_kl\_delta = 0.02} and a second-order trust-region step. The roster is intentionally narrow: the methods are close enough for a matched comparison, but different enough to test increasing levels of policy-step control.

\subsection{Implementation and validation}

The implementation is in the submitted repository under \texttt{src/}, \texttt{configs/}, \texttt{scripts/}, and \texttt{tests/}. The algorithm runners are separated by method: \texttt{a2c\_runner.py}, \texttt{ppo\_clip\_runner.py}, \texttt{ppo\_kl\_runner.py}, and \texttt{trpo\_runner.py}. The PPO-KL runner contains the main custom behavior: it asserts that ratio clipping is disabled, uses the plain ratio surrogate, and relies on target-KL early stopping. TRPO is library-backed through SB3-Contrib to avoid turning the project into a fragile reimplementation of conjugate-gradient trust-region optimization. That choice is consistent with the evaluation goal: the project is about update-control behavior under a fair harness, not about proving that a student implementation of TRPO can match a mature library.

The codebase includes tests for configuration overrides, environment shapes, logging schema, checkpoint resume behavior, evaluation determinism, normalization, algorithm fairness, smoke execution, runtime status accounting, and the validation-probe framework. Those tests are important because small mismatches can change a deep RL comparison. For example, the tests check that the algorithm runners keep matched actor-critic architecture assumptions and that the logging/manifest files have the expected structure. The final project therefore includes both the benchmark code and the guardrails needed to make the comparison inspectable.

\section{Evaluation}

All quantitative values in this section were recomputed directly from per-run \texttt{metrics.csv}, \texttt{updates.csv}, and \texttt{run\_status.json} files because the aggregate analysis pipeline produced inconsistent outputs.

\begin{table}[t]
  \caption{Mean final evaluation return by environment (mean $\pm$ sd across 5 seeds), with collapsed runs in parentheses. Final return means the last logged evaluation checkpoint for each run, and collapsed runs remain included at their last logged evaluation return.}
  \label{tab:main-results}
  \centering
  \scriptsize
  \resizebox{\linewidth}{!}{%
  \begin{tabular}{lrrrr}
    \toprule
    Environment & A2C & PPO-Clip & PPO-KL & TRPO \\
    \midrule
    HalfCheetah-v4 & $-1.4 \pm 1.7$ (3/5) & $2243.6 \pm 506.4$ (0/5) & $1852.3 \pm 450.9$ (0/5) & $\mathbf{3578.7 \pm 1053.2}$ (0/5) \\
    Hopper-v4      & $243.1 \pm 44.3$ (0/5) & $560.7 \pm 33.0$ (0/5) & $808.0 \pm 83.0$ (0/5) & $\mathbf{1912.1 \pm 1215.8}$ (0/5) \\
    Pendulum-v1    & $-1185.7 \pm 116.7$ (0/5) & $-1206.7 \pm 120.3$ (1/5) & $-1110.3 \pm 137.8$ (0/5) & $\mathbf{-782.0 \pm 118.9}$ (0/5) \\
    Walker2d-v4    & $321.6 \pm 17.1$ (0/5) & $675.1 \pm 84.6$ (0/5) & $504.9 \pm 290.0$ (2/5) & $\mathbf{3256.2 \pm 576.8}$ (0/5) \\
    \bottomrule
  \end{tabular}
  }
\end{table}

Across the 80 planned runs, 74 reached the target budget and 6 terminated early under the collapse rule. Collapse counts were A2C \texttt{3}, PPO-Clip \texttt{1}, PPO-KL \texttt{2}, and TRPO \texttt{0}. Wall-clock cost preserved the expected tradeoff but did not reverse the ranking: A2C averaged \texttt{132.3 s}, PPO-KL \texttt{167.0 s}, TRPO \texttt{249.5 s}, and PPO-Clip \texttt{295.5 s}. These timings are hardware-dependent, but they still matter for an honest comparison.

\begin{table}[t]
  \caption{Run-level completion and collapse summary across the 80 planned benchmark runs.}
  \label{tab:collapse}
  \centering
  \small
  \begin{tabular}{lrrr}
    \toprule
    Algorithm & Planned runs & Collapsed runs & Completed to target budget \\
    \midrule
    A2C & 20 & 3 & 17 \\
    PPO-Clip & 20 & 1 & 19 \\
    PPO-KL & 20 & 2 & 18 \\
    TRPO & 20 & 0 & 20 \\
    \midrule
    Total & 80 & 6 & 74 \\
    \bottomrule
  \end{tabular}
\end{table}

\subsection{Evaluation procedure and data handling}

The evaluation was designed to avoid three common sources of misleading reinforcement-learning results. First, every planned run is represented in the manifest, and runs that collapsed are counted as collapsed rather than quietly excluded. Collapsed runs therefore remain in the final-return mean at their last logged evaluation value. Second, the reported scores use the periodic evaluation callback with deterministic policy actions and a fixed evaluation cadence, rather than training return affected by exploration noise, normalization state, or partial rollout episodes. Third, the report treats \texttt{metrics.csv}, \texttt{updates.csv}, and \texttt{run\_status.json} as the source of truth because the final audit found inconsistent aggregate exports. This makes the final narrative more manual than ideal, but ties the reported claims to run-level evidence rather than stale summary files.

The evaluation also separates sample-efficiency and compute-cost concerns. The main table compares final evaluation return under matched environment-step budgets, while the wall-clock summary reports the cost of reaching those results on the local setup. TRPO being slower is expected because its update is more complex. The useful question is whether its extra cost buys enough robustness and return to matter. In this benchmark, it did under the default configuration, but the conclusion remains bounded by the limited seed count and missing sensitivity sweeps.

\subsection{Main findings}

The main locomotion tasks support a narrow but defensible conclusion: explicit step control improved default-setting performance and run-level robustness relative to A2C. HalfCheetah-v4 and Walker2d-v4 carry most of that evidence. On HalfCheetah-v4, A2C averaged \texttt{-1.4} with \texttt{3/5} collapses, PPO-Clip and PPO-KL improved to \texttt{2243.6} and \texttt{1852.3} with zero collapses, and TRPO reached \texttt{3578.7}. On Walker2d-v4, TRPO separated even more strongly at \texttt{3256.2}, PPO-Clip reached \texttt{675.1}, PPO-KL reached \texttt{504.9} with \texttt{2/5} collapses, and A2C remained low at \texttt{321.6}.

Hopper-v4 points in the same direction but with much wider uncertainty. TRPO had the highest mean final return (\texttt{1912.1}), but the standard deviation (\texttt{1215.8}) is large enough that the TRPO versus PPO-KL margin there should not be described as statistically clean. Pendulum-v1 shows the same qualitative ordering, but it used a smaller 100k-step budget and mainly serves as a sanity-check environment rather than the main evidence.

TRPO therefore has the strongest overall empirical profile in this benchmark, but that claim should be stated precisely: it had the highest mean final return on all four environments and zero collapses across 20 runs, while Hopper-v4 still carries wide uncertainty. PPO-Clip was the strongest first-order baseline. It consistently improved on A2C in the locomotion suite and suffered only one collapse in the benchmark.

PPO-KL behaved differently from PPO-Clip in implementation terms: on the MuJoCo tasks, policy optimization stopped after only about \texttt{2.7 to 3.1} epochs on average rather than the nominal maximum of \texttt{10}. That shows the early-stop rule was active. It does not, by itself, prove that KL-based stopping caused the observed return differences rather than simply co-occurring with them, so the claim should remain narrow.

\subsection{Interpretation by environment}

HalfCheetah-v4 is the clearest example of the project thesis. A2C failed to make useful progress under the locked default configuration and collapsed in three of five seeds. PPO-Clip and PPO-KL both avoided collapse and achieved much higher final returns. TRPO achieved the highest mean return, but the more important result is that all explicit step-control methods were much more usable than A2C. This supports the claim that controlling policy movement matters when the continuous-control task requires sustained improvement rather than a short initial gain.

Walker2d-v4 is the strongest result for TRPO. The gap between TRPO and the first-order methods was large, and TRPO had no collapses. PPO-KL was less reliable here, with two collapses and a large standard deviation. That result matters because it prevents an oversimplified conclusion that any KL-based control is automatically enough. The details of the control mechanism still matter. In this benchmark, target-KL early stopping was active, but it did not match the stability of the explicit TRPO trust-region update on Walker2d-v4.

Hopper-v4 is more ambiguous. TRPO again had the highest mean final return, but the standard deviation was very large. The correct interpretation is directional rather than definitive: Hopper-v4 is consistent with the broader conclusion that explicit step control helps, but the TRPO versus PPO-KL ranking on this environment should not be treated as statistically clean. Pendulum-v1 is the smallest task and was run with a smaller step budget. It is useful as a sanity check that the implementations train and log properly, but the main empirical weight should be placed on the MuJoCo locomotion tasks.

\subsection{Robustness and metric disagreement}

For this paper's main thesis about robustness, collapse incidence is the primary measure. It is a run-level failure outcome and directly reflects whether training remained usable. By that metric the ranking is clean: A2C was worst, PPO-Clip and PPO-KL were intermediate, and TRPO was best.

The update-level unstable-update metric tells a less tidy story. TRPO logged nonzero unstable-update counts, concentrated on Hopper-v4 and Walker2d-v4, while the other methods did not. That diagnostic should not be over-interpreted. First, it is measuring something local rather than a run-level outcome. Second, it uses a tighter KL threshold for TRPO and PPO-KL (\texttt{0.04}) than for A2C and PPO-Clip (\texttt{0.05}). The correct conclusion is therefore not that TRPO was globally unstable, but that the local instability diagnostic and the run-level collapse metric capture different phenomena. For the main benchmark claim, collapse incidence is the more decision-relevant robustness measure.

This disagreement is also a useful lesson about empirical RL evaluation. A local update can look suspicious under a KL-and-drop diagnostic without causing the whole run to fail, and a run can still be successful even if a few updates are locally noisy. Conversely, a method can avoid a particular update-level flag while still producing poor returns or collapsed runs. That is why the final report treats collapse incidence as the primary robustness metric and update instability as diagnostic context. The project would be weaker if it reported only final return, because collapse behavior is exactly the kind of failure that a stability-focused benchmark should expose.

\subsection{Practical implications}

For a practitioner choosing among these methods under a limited compute budget, the results suggest three practical lessons. First, A2C is useful as a reference point but should not be treated as a reliable default for these continuous-control tasks. It was cheaper per run, but that lower cost did not compensate for poor HalfCheetah-v4 performance and three collapses. Second, PPO-Clip remains a strong engineering compromise: it avoids TRPO's complexity, uses a first-order update, improves clearly on A2C in the locomotion suite, and had only one collapse. Third, PPO-KL shows why method names are not enough. Its unclipped ratio surrogate with target-KL early stopping was active, but it did not guarantee the best robustness profile; on Walker2d-v4, PPO-KL had two collapses and a high standard deviation. The result is not that KL information is useless, but that the way KL is used matters.

The strongest practical case for TRPO is reliability under this matched default setup. TRPO was slower than A2C and PPO-KL, but it had zero collapses and the highest mean final return on all four environments. That does not make TRPO universally best, especially because Hopper-v4 variance was high and the sensitivity sweeps were not completed. It does mean that the extra trust-region machinery bought something measurable in this benchmark. For a stability-focused project, that tradeoff is central: the method with the heavier update was also the one that most consistently avoided failed runs.

The broader implication is that continuous-control policy optimization should be evaluated as a system, not only as an objective function. Architecture, rollout length, normalization, value bootstrapping, evaluation cadence, collapse accounting, and implementation details all affect the apparent algorithm ranking. The project's matched protocol is valuable because it makes those assumptions visible. Even when the benchmark does not produce a universal answer, it gives a more trustworthy answer to the question it actually asks.

\subsection{Threats to validity}

The largest limitation is the missing hyperparameter-sensitivity sweep analysis. The benchmark therefore supports a default-setting claim much more strongly than a robustness-to-hyperparameter-variation claim. The original project plan included sensitivity sweeps over A2C learning rate, PPO-Clip learning rate and clip range, PPO-KL learning rate and target KL, and TRPO trust-region radius. Those sweeps would have made it possible to say whether one method remains stable across a broader tuning range. Without them, the project should not claim universal robustness.

The second limitation is the aggregate analysis bug: report claims had to be rebuilt from raw per-run artifacts. This is not ideal, but the mitigation is transparent. The final paper uses per-run \texttt{metrics.csv}, \texttt{updates.csv}, and \texttt{run\_status.json} files as the source of truth and explicitly warns against blindly citing known-bad aggregate exports. Third, five seeds per condition are reasonable for a one-student project but still produce wide uncertainty on noisy tasks such as Hopper-v4. Fourth, the benchmark is intentionally small. Four environments are enough for a disciplined course comparison, but not enough to justify universal claims about policy optimization. Finally, the results remain implementation-dependent. The project intentionally uses a library-backed, fairness-controlled setup, and that is the right engineering choice here, but it limits how abstractly the ranking should be interpreted.

\subsection{Reproducibility and submission artifacts}

The submitted repository is \url{https://github.com/Hardik-S/RL-Project}. The intended repository root is \texttt{RL-Project}, not the broader workspace folder that contains drafts, lecture files, presentation materials, and other local artifacts. The repository contains the training code, configuration files, tests, documentation, manifests, and final report artifacts. The reproducibility notes in \texttt{README.md} and \texttt{docs/REPRODUCIBILITY.md} define the setup and commands needed to run tests, smoke checks, benchmark jobs, and analysis. The project also records planned and completed runs so that failed or collapsed jobs are visible rather than silently omitted.

\section{Conclusion}

This project asked whether explicit control of policy-update size improves continuous-control training behavior under a matched benchmark. Under the locked default configuration, the answer is yes relative to A2C on the main locomotion tasks. TRPO produced the strongest overall empirical profile, combining the highest mean final return with zero collapses, although Hopper-v4 variance means that claim should be stated with caution rather than as a blanket dominance result. PPO-Clip was the strongest simpler first-order baseline. PPO-KL was active and competitive, but less reliable on Walker2d-v4.

The developed technique, a matched update-stability benchmark with explicit collapse accounting, can effectively tackle the evaluation problem posed in this project: it separates update-control behavior from several avoidable implementation confounds and makes failed runs visible. It does not prove that one optimizer is always best, and it does not replace a full hyperparameter sweep. The most defensible summary is therefore modest: in this benchmark, stronger step control helped, and explicit trust-region updates worked best under the tested defaults.

Future research should first repair and harden the aggregate analysis pipeline so that report tables are generated automatically from the same per-run source of truth used in this paper. The next step should be the missing sensitivity sweeps, because a stability claim is much stronger if it holds across learning rates, clip ranges, target-KL thresholds, and trust-region radii. A larger follow-up would add more continuous-control tasks, increase the seed count, report confidence intervals or paired statistical comparisons, and compare target-KL early stopping against an adaptive KL-penalty variant.

\clearpage
\section*{References}

{\small

[1] Ronald J.~Williams. Simple Statistical Gradient-Following Algorithms for Connectionist Reinforcement Learning. \emph{Machine Learning}, 8(3--4):229--256, 1992.

[2] Volodymyr Mnih, Adria Puigdomenech Badia, Mehdi Mirza, Alex Graves, Timothy P.~Lillicrap, Tim Harley, David Silver, and Koray Kavukcuoglu. Asynchronous Methods for Deep Reinforcement Learning. In \emph{Proceedings of the 33rd International Conference on Machine Learning}, 2016.

[3] Sham Kakade and John Langford. Approximately Optimal Approximate Reinforcement Learning. In \emph{Proceedings of the Nineteenth International Conference on Machine Learning}, 2002.

[4] John Schulman, Sergey Levine, Pieter Abbeel, Michael I.~Jordan, and Philipp Moritz. Trust Region Policy Optimization. In \emph{Proceedings of the 32nd International Conference on Machine Learning}, 2015.

[5] John Schulman, Philipp Moritz, Sergey Levine, Michael I.~Jordan, and Pieter Abbeel. High-Dimensional Continuous Control Using Generalized Advantage Estimation. arXiv:1506.02438, 2015.

[6] John Schulman, Filip Wolski, Prafulla Dhariwal, Alec Radford, and Oleg Klimov. Proximal Policy Optimization Algorithms. arXiv:1707.06347, 2017.

[7] Peter Henderson, Riashat Islam, Philip Bachman, Joelle Pineau, Doina Precup, and David Meger. Deep Reinforcement Learning That Matters. In \emph{Proceedings of the AAAI Conference on Artificial Intelligence}, 2018.

[8] Logan Engstrom, Andrew Ilyas, Shibani Santurkar, Dimitris Tsipras, Firdaus Janoos, Larry Rudolph, and Aleksander Madry. Implementation Matters in Deep Policy Gradients: A Case Study on PPO and TRPO. In \emph{International Conference on Learning Representations}, 2020.

}

\end{document}
